\chapter{EXPERIMENTAL RESULTS AND DISCUSSION}
\label{ch:experimental_results_and_discussion}

\hs This chapter presents the empirical results obtained from the experimental protocol described in Section~\ref{sec:experimental_design} of Chapter~\ref{ch:methodology}. The results are reported separately for classification and regression tasks. Each table summarizes the mean and standard deviation across five repeated train--test splits using the random seeds defined in the experimental design.

The objective of this chapter is to report and discuss the observed model behavior without modifying the numerical outputs produced by the experiment. The evaluation includes baseline supervised learning models, COBRA-based aggregation models, and \KFCProcedure{}-based models. Classification experiments compare standard classifiers with \varname{CombinedClassifier} and \varname{KFCClassifier}. Regression experiments compare standard regressors with \varname{GradientCOBRA}, \varname{MixCOBRARegressor}, and \varname{KFCRegressor} configured with a \varname{GradientCOBRA} combiner.

\section{Classification Results}
\label{sec:classification_results}

\hs The classification results are reported as mean $\pm$ standard deviation over five random seeds. The evaluated datasets include Heart Disease, Diabetes, Loan Prediction, and Cyber Security Attacks. The reported metrics are Accuracy, Precision, Recall, and F1-score.

\subsection{Heart Disease Classification}

\begin{table}[H]
    \centering
    \caption{Mean Heart Disease Classification Results}
    \label{tab:heart_aggregated_results}
    \small
    \resizebox{\textwidth}{!}{%
        \begin{tabular}{lcccc}
            \toprule
            \textbf{Model}                   & \textbf{Accuracy} & \textbf{Precision} & \textbf{Recall}   & \textbf{F1 Score} \\
            \midrule
            KFCClassifier-CombinedClassifier & $0.819 \pm 0.037$ & $0.821 \pm 0.038$  & $0.819 \pm 0.037$ & $0.819 \pm 0.036$ \\
            Logistic Regression              & $0.806 \pm 0.032$ & $0.810 \pm 0.034$  & $0.806 \pm 0.032$ & $0.805 \pm 0.033$ \\
            Random Forest Classifier         & $0.806 \pm 0.060$ & $0.811 \pm 0.065$  & $0.806 \pm 0.060$ & $0.805 \pm 0.060$ \\
            Linear SVC                       & $0.800 \pm 0.042$ & $0.805 \pm 0.042$  & $0.800 \pm 0.042$ & $0.798 \pm 0.044$ \\
            KNeighborsClassifier             & $0.787 \pm 0.108$ & $0.788 \pm 0.108$  & $0.787 \pm 0.108$ & $0.787 \pm 0.108$ \\
            CombinedClassifier               & $0.768 \pm 0.062$ & $0.779 \pm 0.068$  & $0.768 \pm 0.062$ & $0.764 \pm 0.061$ \\
            GaussianNB                       & $0.723 \pm 0.093$ & $0.743 \pm 0.075$  & $0.723 \pm 0.093$ & $0.717 \pm 0.098$ \\
            \bottomrule
        \end{tabular}%
    }
\end{table}



\hs Table~\ref{tab:heart_aggregated_results} presents the mean classification results for the Heart Disease dataset across five repeated train--test splits. Because this dataset contains only about 302 de-duplicated observations and the experiment uses a test size of 0.10, each split evaluates roughly 30 test samples. Under this small-test-set setting, \textbf{KFCClassifier-CombinedClassifier} obtains the nominally highest mean Accuracy and F1-score, both equal to $0.819$. However, its lead over Logistic Regression and Random Forest Classifier, which both obtain a mean Accuracy of $0.806$ and an F1-score of $0.805$, is small and lies within the reported standard deviations. Therefore, this result should be interpreted as an encouraging descriptive result rather than statistically reliable evidence of superiority.





Compared with the standalone \textbf{CombinedClassifier}, the KFC-based model increases the mean Accuracy from $0.768$ to $0.819$. This suggests that the clustering and local learning stages may provide more informative prediction structures for this dataset. Nevertheless, because the test set is small and the standard deviations overlap with strong baselines, the Heart Disease result is used only as descriptive evidence that the KFCProcedure workflow is functional and potentially useful, not as conclusive evidence that it is superior.

Overall, the Heart Disease result should be interpreted cautiously. It shows that the KFC classification workflow can be executed successfully and can produce competitive descriptive performance on this dataset, but it does not provide statistically confirmed evidence of superiority over the strongest baseline classifiers.

\subsection{Diabetes Classification}

\begin{table}[H]
    \centering
    \caption{Mean Diabetes Classification Results}
    \label{tab:diabetes_aggregated_results}
    \small
    \resizebox{\textwidth}{!}{%
        \begin{tabular}{lcccc}
            \toprule
            \textbf{Model}                   & \textbf{Accuracy} & \textbf{Precision} & \textbf{Recall}   & \textbf{F1 Score} \\
            \midrule
            Random Forest Classifier         & $0.761 \pm 0.036$ & $0.755 \pm 0.041$  & $0.761 \pm 0.036$ & $0.754 \pm 0.042$ \\
            KFCClassifier-CombinedClassifier & $0.757 \pm 0.056$ & $0.755 \pm 0.066$  & $0.757 \pm 0.056$ & $0.741 \pm 0.063$ \\
            Logistic Regression              & $0.755 \pm 0.047$ & $0.748 \pm 0.051$  & $0.755 \pm 0.047$ & $0.744 \pm 0.052$ \\
            Linear SVC                       & $0.753 \pm 0.047$ & $0.747 \pm 0.051$  & $0.753 \pm 0.047$ & $0.742 \pm 0.052$ \\
            CombinedClassifier               & $0.751 \pm 0.050$ & $0.745 \pm 0.060$  & $0.751 \pm 0.050$ & $0.734 \pm 0.061$ \\
            GaussianNB                       & $0.740 \pm 0.040$ & $0.737 \pm 0.040$  & $0.740 \pm 0.040$ & $0.737 \pm 0.041$ \\
            KNeighborsClassifier             & $0.734 \pm 0.047$ & $0.726 \pm 0.053$  & $0.734 \pm 0.047$ & $0.727 \pm 0.052$ \\
            \bottomrule
        \end{tabular}%
    }
\end{table}

\hs Table~\ref{tab:diabetes_aggregated_results} presents the mean classification results for the Diabetes dataset. Random Forest Classifier achieves the highest reported mean Accuracy and F1-score, with values of $0.761$ and $0.754$, respectively. \textbf{KFCClassifier-CombinedClassifier} obtains the second-highest mean Accuracy of $0.757$ and the same mean Precision as Random Forest Classifier, at $0.755$. This shows that the KFC classification configuration remains competitive on this dataset and performs close to the strongest baseline model in terms of overall classification accuracy.

Compared with the standalone \textbf{CombinedClassifier}, the KFC-based model improves the mean Accuracy from $0.751$ to $0.757$ and the F1-score from $0.734$ to $0.741$. This suggests that the clustering and local learning stages may provide some additional benefit before the final classification aggregation step. However, the improvement is modest, and the F1-score of the KFC model remains slightly lower than Random Forest Classifier and Logistic Regression.

Therefore, the Diabetes results suggest that KFCProcedure is functional and competitive, but the latent cluster structure in this dataset may not be strong enough to clearly outperform the best global baseline model. Since the standard deviations overlap across several models, these differences should be interpreted descriptively rather than as statistically confirmed performance differences.

\subsection{Loan Prediction Classification}

\begin{table}[H]
    \centering
    \caption{Mean Loan Prediction Classification Results}
    \label{tab:loan_prediction_aggregated_results}
    \small
    \resizebox{\textwidth}{!}{%
        \begin{tabular}{lcccc}
            \toprule
            \textbf{Model}                   & \textbf{Accuracy} & \textbf{Precision} & \textbf{Recall}   & \textbf{F1 Score} \\
            \midrule
            Logistic Regression              & $0.826 \pm 0.024$ & $0.851 \pm 0.020$  & $0.826 \pm 0.024$ & $0.804 \pm 0.032$ \\
            Linear SVC                       & $0.824 \pm 0.024$ & $0.850 \pm 0.020$  & $0.824 \pm 0.024$ & $0.801 \pm 0.031$ \\
            CombinedClassifier               & $0.821 \pm 0.028$ & $0.841 \pm 0.029$  & $0.821 \pm 0.028$ & $0.799 \pm 0.034$ \\
            KFCClassifier-CombinedClassifier & $0.820 \pm 0.030$ & $0.837 \pm 0.031$  & $0.820 \pm 0.030$ & $0.798 \pm 0.037$ \\
            GaussianNB                       & $0.818 \pm 0.025$ & $0.828 \pm 0.028$  & $0.818 \pm 0.025$ & $0.799 \pm 0.029$ \\
            Random Forest Classifier         & $0.795 \pm 0.023$ & $0.795 \pm 0.034$  & $0.795 \pm 0.023$ & $0.780 \pm 0.024$ \\
            KNeighborsClassifier             & $0.777 \pm 0.036$ & $0.783 \pm 0.050$  & $0.777 \pm 0.036$ & $0.748 \pm 0.042$ \\
            \bottomrule
        \end{tabular}%
    }
\end{table}




\hs Table~\ref{tab:loan_prediction_aggregated_results} presents the mean classification results for the Loan Prediction dataset. Logistic Regression achieves the highest reported mean performance, with Accuracy of $0.826$, Precision of $0.851$, and F1-score of $0.804$. Linear SVC performs very closely, with Accuracy of $0.824$ and F1-score of $0.801$. These results indicate that linear models are highly effective for this dataset after preprocessing, suggesting that the relationship between the selected features and the target variable may be captured reasonably well by relatively simple decision boundaries.

The aggregation-based models also remain competitive. \varname{CombinedClassifier} obtains a mean Accuracy of $0.821$, while \varname{KFCClassifier-CombinedClassifier} obtains a mean Accuracy of $0.820$. These values are close to the best-performing linear baselines, but they do not exceed them. Compared with the standalone \varname{CombinedClassifier}, the KFC-based model shows a very similar result, indicating that the clustering and local learning stages do not provide a clear additional advantage for this dataset under the tested configuration.

Overall, the Loan Prediction results suggest that KFCProcedure is functional and comparable to strong baseline models, but the dataset may not contain sufficiently strong latent cluster structures for the clusterwise approach to clearly improve performance. Since the differences among the top models are small and their standard deviations overlap, the results should be interpreted descriptively rather than as statistically confirmed superiority of one model over another.

\subsection{Cyber Security Attacks Classification}

\begin{table}[H]
    \centering
    \caption{Mean Cyber Security Attacks Classification Results}
    \label{tab:cyber_security_aggregated_results}
    \small
    \resizebox{\textwidth}{!}{%
        \begin{tabular}{lcccc}
            \toprule
            \textbf{Model}                   & \textbf{Accuracy} & \textbf{Precision} & \textbf{Recall}   & \textbf{F1 Score} \\
            \midrule
            Linear SVC                       & $0.340 \pm 0.007$ & $0.340 \pm 0.007$  & $0.340 \pm 0.007$ & $0.339 \pm 0.007$ \\
            Logistic Regression              & $0.340 \pm 0.006$ & $0.339 \pm 0.006$  & $0.340 \pm 0.006$ & $0.338 \pm 0.007$ \\
            KFCClassifier-CombinedClassifier & $0.339 \pm 0.013$ & $0.339 \pm 0.014$  & $0.339 \pm 0.013$ & $0.327 \pm 0.010$ \\
            KNeighborsClassifier             & $0.335 \pm 0.008$ & $0.337 \pm 0.011$  & $0.335 \pm 0.008$ & $0.328 \pm 0.009$ \\
            CombinedClassifier               & $0.334 \pm 0.004$ & $0.288 \pm 0.064$  & $0.334 \pm 0.004$ & $0.292 \pm 0.047$ \\
            Random Forest Classifier         & $0.333 \pm 0.012$ & $0.333 \pm 0.012$  & $0.333 \pm 0.012$ & $0.333 \pm 0.012$ \\
            GaussianNB                       & $0.331 \pm 0.011$ & $0.330 \pm 0.011$  & $0.331 \pm 0.011$ & $0.329 \pm 0.011$ \\
            \bottomrule
        \end{tabular}%
    }
\end{table}




\hs Table~\ref{tab:cyber_security_aggregated_results} presents the mean classification results for the Cyber Security Attacks dataset. All evaluated models obtain very close Accuracy values, ranging from approximately $0.331$ to $0.340$. Since the target has three classes, a simple random-guessing baseline would be approximately $0.333$. Therefore, the reported results are close to chance level, indicating that the selected features and preprocessing strategy did not extract a strong usable signal for this classification task under the current experimental setting.







Compared with the standalone \varname{CombinedClassifier}, the KFC-based model changes the mean Accuracy from $0.334$ to $0.339$ and the F1-score from $0.292$ to $0.327$. However, because all models remain near the three-class chance baseline, this difference is not substantively meaningful. The Cyber Security Attacks experiment should therefore be interpreted only as evidence that the KFC-based pipeline is executable on a multi-class dataset, not as evidence that it provides a clear performance advantage.

Overall, the Cyber Security Attacks results suggest that this dataset is a difficult classification case under the selected preprocessing and feature representation. The small performance differences among models indicate that the current features may not provide strong separation between classes. Therefore, the result does not provide strong evidence of a clear latent cluster structure that benefits the KFCProcedure approach. Additional feature engineering, more specialized cyber security variables, or a different sampling strategy may be required to improve predictive performance on this dataset.

\section{Regression Results}
\label{sec:regression_results}

\hs The regression results are reported as mean $\pm$ standard deviation over five random seeds. The evaluated datasets include Abalone, California Housing Price, and Wind Turbine. The reported metrics are MAPE, RMSE, $R^2$, and the average value of the true target in the test split. For the Wind Turbine dataset, MAPE is not reported because the target contains zero or near-zero values.

\subsection{Abalone Regression}

\begin{table}[H]
    \centering
    \caption{Mean Abalone Regression Results}
    \label{tab:abalone_aggregated_results}
    \small
    \resizebox{\textwidth}{!}{%
        \begin{tabular}{lcccc}
            \toprule
            \textbf{Model}              & \textbf{MAPE}     & \textbf{RMSE}     & \textbf{$R^2$}    & \textbf{AVG\_Y\_TRUE} \\
            \midrule
            KFCRegressor-GradientCOBRA  & $0.156 \pm 0.005$ & $2.143 \pm 0.056$ & $0.560 \pm 0.012$ & $9.932 \pm 0.058$   \\
            MixCOBRARegressor           & $0.151 \pm 0.006$ & $2.156 \pm 0.064$ & $0.555 \pm 0.018$ & $9.932 \pm 0.058$   \\
            SVR                         & $0.141 \pm 0.006$ & $2.164 \pm 0.066$ & $0.551 \pm 0.017$ & $9.932 \pm 0.058$   \\
            GradientCOBRA               & $0.154 \pm 0.007$ & $2.167 \pm 0.067$ & $0.550 \pm 0.019$ & $9.932 \pm 0.058$   \\
            Random Forest Regressor     & $0.153 \pm 0.005$ & $2.174 \pm 0.075$ & $0.547 \pm 0.021$ & $9.932 \pm 0.058$   \\
            Gradient Boosting Regressor & $0.154 \pm 0.005$ & $2.198 \pm 0.069$ & $0.537 \pm 0.020$ & $9.932 \pm 0.058$   \\
            Linear Regression           & $0.162 \pm 0.005$ & $2.206 \pm 0.047$ & $0.533 \pm 0.014$ & $9.932 \pm 0.058$   \\
            Ridge                       & $0.162 \pm 0.005$ & $2.206 \pm 0.047$ & $0.533 \pm 0.014$ & $9.932 \pm 0.058$   \\
            KNeighborsRegressor         & $0.160 \pm 0.004$ & $2.303 \pm 0.069$ & $0.492 \pm 0.024$ & $9.932 \pm 0.058$   \\
            Lasso                       & $0.214 \pm 0.007$ & $2.693 \pm 0.075$ & $0.305 \pm 0.017$ & $9.932 \pm 0.058$   \\
            \bottomrule
        \end{tabular}%
    }
\end{table}



\hs Table~\ref{tab:abalone_aggregated_results} presents the mean regression results for the Abalone dataset. \varname{KFCRegressor-GradientCOBRA} obtains the nominally highest reported mean $R^2$ value of $0.560$ and the lowest RMSE value of $2.143$ among the evaluated models. \varname{MixCOBRARegressor}, \varname{SVR}, and \varname{GradientCOBRA} are very close, with mean $R^2$ values of $0.555$, $0.551$, and $0.550$, respectively. The differences among these top models fall within approximately one reported standard deviation, and no statistical significance test was conducted. In addition, \varname{SVR} achieves the lowest MAPE value. Therefore, the Abalone result should be described as a nominal advantage for KFCRegressor-GradientCOBRA in RMSE and $R^2$, rather than as a statistically confirmed best overall regression fit.


This result suggests that the Abalone dataset may contain meaningful latent cluster structure, where different regions of the input feature space benefit from local predictive modeling before final aggregation. In this case, the KFCProcedure assumption may be relevant because the clusterwise model obtains a nominal RMSE and $R^2$ advantage over the standalone global models and standalone COBRA-based aggregation models. However, because the differences are small and no statistical significance test was conducted, this should be interpreted as descriptive evidence under the tested configuration rather than as a formal proof of latent clusters.

\subsection{California Housing Price Regression}

\begin{table}[H]
    \centering
    \caption{Mean California Housing Price Regression Results}
    \label{tab:housing_price_aggregated_results}
    \small
    \resizebox{\textwidth}{!}{%
        \begin{tabular}{lcccc}
            \toprule
            \textbf{Model}              & \textbf{MAPE}     & \textbf{RMSE}     & \textbf{$R^2$}     & \textbf{AVG\_Y\_TRUE} \\
            \midrule
            Random Forest Regressor     & $0.012 \pm 0.000$ & $0.213 \pm 0.006$ & $0.874 \pm 0.007$  & $11.906 \pm 0.012$    \\
            Gradient Boosting Regressor & $0.013 \pm 0.000$ & $0.226 \pm 0.004$ & $0.858 \pm 0.008$  & $11.906 \pm 0.012$    \\
            SVR                         & $0.014 \pm 0.001$ & $0.239 \pm 0.011$ & $0.841 \pm 0.015$  & $11.906 \pm 0.012$    \\
            KNeighborsRegressor         & $0.016 \pm 0.000$ & $0.259 \pm 0.009$ & $0.813 \pm 0.012$  & $11.906 \pm 0.012$    \\
            MixCOBRARegressor           & $0.016 \pm 0.003$ & $0.262 \pm 0.040$ & $0.805 \pm 0.064$  & $11.906 \pm 0.012$    \\
            GradientCOBRA               & $0.017 \pm 0.003$ & $0.272 \pm 0.033$ & $0.791 \pm 0.054$  & $11.906 \pm 0.012$    \\
            KFCRegressor-GradientCOBRA  & $0.018 \pm 0.000$ & $0.283 \pm 0.007$ & $0.778 \pm 0.009$  & $11.906 \pm 0.012$    \\
            Lasso                       & $0.042 \pm 0.001$ & $0.600 \pm 0.007$ & $-0.000 \pm 0.001$ & $11.906 \pm 0.012$    \\
            Linear Regression (unstable) & $0.020 \pm 0.002$ & $0.713 \pm 0.576$ & $-1.149 \pm 2.726$ & $11.906 \pm 0.012$    \\
            Ridge (unstable)             & $0.020 \pm 0.002$ & $0.714 \pm 0.577$ & $-1.158 \pm 2.738$ & $11.906 \pm 0.012$    \\
            \bottomrule
        \end{tabular}%
    }
\end{table}

\hs Table~\ref{tab:housing_price_aggregated_results} presents the mean regression results for the California Housing Price dataset, the target variable is log-transformed before modeling. Therefore, the RMSE and MAPE values in Table~\ref{tab:housing_price_aggregated_results} are reported on the logarithmic target scale rather than on the original house-price scale. Random Forest Regressor achieves the strongest overall performance, with the lowest RMSE value of $0.213$ and the highest mean $R^2$ value of $0.874$. Gradient Boosting Regressor and SVR also perform strongly, with mean $R^2$ values of $0.858$ and $0.841$, respectively.

The Linear Regression and Ridge rows show negative mean $R^2$ values with very large standard deviations, which indicates numerical instability across the five random seeds rather than a reliable substantive result. A negative $R^2$ means that the model performs worse than predicting the test-set mean. This instability may be related to the deterministic first-5,000-row sampling strategy, possible ordering effects in the dataset, feature collinearity, or insufficient regularization and scaling for linear models. Therefore, these two rows are retained only as diagnostic evidence of an experimental limitation and should not be used to draw conclusions about the true capability of linear regression on the California Housing dataset.

The COBRA-based models remain executable and competitive, but they do not outperform the strongest tree-based baseline models on this dataset. \varname{MixCOBRARegressor} obtains a mean $R^2$ of $0.805$, while \varname{GradientCOBRA} obtains $0.791$. \varname{KFCRegressor-GradientCOBRA} obtains a lower mean $R^2$ of $0.778$ and a higher RMSE of $0.283$ compared with the best baseline models.

These results suggest that, under the tested configuration, the California Housing Price dataset is better captured by strong global nonlinear models such as Random Forest and Gradient Boosting. The clusterwise KFC configuration remains functional, but it does not provide a clear performance advantage. This may indicate that the chosen clustering structure, number of clusters, or local model configuration does not fully match the underlying structure of this dataset. Therefore, the result does not provide strong evidence that the KFCProcedure latent-cluster assumption is beneficial for this dataset under the current experimental setting.

\subsection{Wind Turbine Regression}

\begin{table}[H]
    \centering
    \caption{Mean Wind Turbine Regression Results}
    \label{tab:wind_turbine_aggregated_results}
    \small
    \resizebox{\textwidth}{!}{%
        \begin{tabular}{lcccc}
            \toprule
            \textbf{Model}              & \textbf{MAPE}                                     & \textbf{RMSE}        & \textbf{$R^2$}    & \textbf{AVG\_Y\_TRUE} \\
            \midrule
            MixCOBRARegressor          & -- & $519.079 \pm 19.711$ & $0.876 \pm 0.009$ & $1591.707 \pm 43.202$ \\
            Gradient Boosting Regressor & -- & $519.241 \pm 17.145$ & $0.876 \pm 0.008$ & $1591.707 \pm 43.202$ \\
            GradientCOBRA              & -- & $521.699 \pm 20.311$ & $0.874 \pm 0.009$ & $1591.707 \pm 43.202$ \\
            Random Forest Regressor    & -- & $526.608 \pm 10.205$ & $0.872 \pm 0.005$ & $1591.707 \pm 43.202$ \\
            KNeighborsRegressor        & -- & $538.703 \pm 23.032$ & $0.866 \pm 0.011$ & $1591.707 \pm 43.202$ \\
            KFCRegressor-GradientCOBRA & -- & $583.174 \pm 23.157$ & $0.843 \pm 0.013$ & $1591.707 \pm 43.202$ \\
            Lasso                      & -- & $740.766 \pm 15.774$ & $0.747 \pm 0.011$ & $1591.707 \pm 43.202$ \\
            Linear Regression          & -- & $740.780 \pm 15.814$ & $0.747 \pm 0.011$ & $1591.707 \pm 43.202$ \\
            Ridge                      & -- & $740.785 \pm 15.805$ & $0.747 \pm 0.011$ & $1591.707 \pm 43.202$ \\
            SVR                        & -- & $939.607 \pm 15.741$ & $0.593 \pm 0.011$ & $1591.707 \pm 43.202$ \\
            \bottomrule
        \end{tabular}%
    }

    \vspace{0.2cm}
\begin{minipage}{0.95\textwidth}
    \footnotesize \textit{Note:} MAPE is not reported for this dataset because the Wind Turbine target contains zero or near-zero values, which make percentage errors undefined or numerically unstable. The discussion therefore relies on RMSE and $R^2$.
\end{minipage}
\end{table}




\hs Table~\ref{tab:wind_turbine_aggregated_results} presents the mean regression results for the Wind Turbine dataset. \varname{MixCOBRARegressor} and Gradient Boosting Regressor obtain the strongest reported performance, with very similar mean $R^2$ values of $0.876$. Their RMSE values are also almost identical, with \varname{MixCOBRARegressor} achieving $519.079$ and Gradient Boosting Regressor achieving $519.241$. \varname{GradientCOBRA} and Random Forest Regressor also perform competitively, with mean $R^2$ values of $0.874$ and $0.872$, respectively.




MAPE is not reported in Table~\ref{tab:wind_turbine_aggregated_results} because the target variable contains values equal to or very close to zero. Since MAPE divides the prediction error by the true target value, very small denominators can produce undefined or inflated percentage errors. Therefore, the Wind Turbine discussion relies mainly on RMSE and $R^2$.

\varname{KFCRegressor-GradientCOBRA} completes successfully, with a mean $R^2$ value of $0.843$ and an RMSE of $583.174$. Although this confirms that the KFC regression pipeline is executable on this dataset, its performance is lower than the strongest standalone COBRA-based and tree-based models. This suggests that, under the tested configuration, the selected clusterwise structure and local linear models do not provide a clear advantage for the Wind Turbine dataset. The result may indicate that the relationship between wind speed, theoretical power curve, wind direction, and active power is better captured by strong global nonlinear models or direct aggregation methods than by the current KFC configuration.

\section{Overall Discussion}
\label{sec:overall_discussion_results}

\hs The experimental results show that the implemented \textbf{kfc\_procedure} package can support a unified evaluation workflow across heterogeneous supervised learning tasks. The framework is executable for both classification and regression experiments, and it can be compared with standard baseline models and COBRA-based aggregation models using repeated train--test evaluations. Across the seven datasets, the results demonstrate that the proposed framework is functional, reusable, and comparable, although its predictive advantage depends on the characteristics of each dataset.

% For classification, \varname{KFCClassifier} is successfully executed on all four classification datasets. It obtains the highest reported mean performance on the Heart Disease dataset, where it slightly outperforms Logistic Regression and Random Forest Classifier. This result suggests that the Heart Disease dataset may contain local patterns or latent subgroups that benefit from the clusterwise learning workflow before final aggregation. On the Diabetes dataset, \varname{KFCClassifier} remains competitive and performs close to Random Forest Classifier, but it does not clearly exceed the strongest baseline. On the Loan Prediction dataset, Logistic Regression and Linear SVC perform slightly better, suggesting that this dataset may be well captured by simpler linear decision boundaries after preprocessing. On the Cyber Security Attacks dataset, all models produce very similar Accuracy values, indicating that the selected features and preprocessing strategy may not provide strong class separation. Overall, the classification results support the functional validity of the KFC classification pipeline, but they do not establish universal superiority over standard classifiers.



For classification, \varname{KFCClassifier} is successfully executed on all four classification datasets. It obtains the nominally highest reported mean performance on the Heart Disease dataset, but the result is based on a small test set and overlaps with strong baselines within the reported standard deviations. On the Diabetes dataset, \varname{KFCClassifier} remains close to Random Forest Classifier, but it does not clearly exceed the strongest baseline. On the Loan Prediction dataset, Logistic Regression and Linear SVC perform slightly better, suggesting that this dataset may be well captured by simpler linear decision boundaries after preprocessing. On the Cyber Security Attacks dataset, all models perform near the three-class chance baseline, indicating that the selected features and preprocessing strategy may not provide strong class separation. Overall, the classification results support the functional validity of the KFC classification pipeline, but they do not establish universal superiority over standard classifiers.

% For regression, \varname{KFCRegressor} with \varname{GradientCOBRA} is successfully executed on all three regression datasets. On the Abalone dataset, \varname{KFCRegressor-GradientCOBRA} achieves the highest reported mean $R^2$ and the lowest RMSE among the evaluated models. This result provides empirical support for the assumption that the Abalone dataset may contain meaningful latent cluster structure, where local predictive modeling can improve performance before final aggregation. However, on the California Housing Price dataset, the strongest performance is obtained by tree-based models, especially Random Forest Regressor and Gradient Boosting Regressor. Similarly, on the Wind Turbine dataset, \varname{MixCOBRARegressor}, Gradient Boosting Regressor, \varname{GradientCOBRA}, and Random Forest Regressor perform better than \varname{KFCRegressor-GradientCOBRA}. These results indicate that the KFC regression pipeline is operational, but its advantage is dataset-dependent and strongly influenced by the suitability of the clustering structure, the selected number of clusters, and the local model configuration.


For regression, \varname{KFCRegressor} with \varname{GradientCOBRA} is successfully executed on all three regression datasets. On the Abalone dataset, \varname{KFCRegressor-GradientCOBRA} obtains a nominal advantage in RMSE and $R^2$, but the differences among the top models are small and fall within approximately one reported standard deviation. Therefore, this result should be interpreted as encouraging descriptive evidence rather than statistically confirmed superiority. On the California Housing Price dataset, the strongest performance is obtained by tree-based models, especially Random Forest Regressor and Gradient Boosting Regressor. Similarly, on the Wind Turbine dataset, \varname{MixCOBRARegressor}, Gradient Boosting Regressor, \varname{GradientCOBRA}, and Random Forest Regressor perform better than \varname{KFCRegressor-GradientCOBRA}. These results indicate that the KFC regression pipeline is operational, but its advantage is dataset-dependent and strongly influenced by the suitability of the clustering structure, the selected number of clusters, and the local model configuration.

The COBRA-based models also show competitive behavior in several cases. \varname{MixCOBRARegressor} performs close to the best model on Abalone and obtains one of the strongest results on the Wind Turbine dataset. \varname{GradientCOBRA} is consistently executable across regression datasets and remains competitive in some cases, especially when prediction-space aggregation is useful. \varname{CombinedClassifier} is also executable for classification, although its predictive performance varies depending on the dataset. The comparison between standalone COBRA-based models and KFC-based models shows that clusterwise learning does not automatically improve performance in every case. Instead, its benefit appears when the dataset contains useful latent group structures and when different regions of the input space benefit from different local predictive models.

Overall, the experimental findings demonstrate that the proposed framework is usable and comparable across multiple supervised learning settings. The results show that the framework can integrate clustering, local predictive modeling, and aggregation-based prediction in a consistent workflow. In particular, the Heart Disease and Abalone results provide encouraging empirical evidence that the KFCProcedure assumption can be relevant for some datasets. However, the results on Diabetes, Loan Prediction, Cyber Security Attacks, California Housing Price, and Wind Turbine show that the advantage of KFCProcedure is not universal. Therefore, the current experiments should be interpreted mainly as a functional and comparative evaluation of the implemented framework, rather than as final evidence that the proposed models always outperform baseline methods.

\section{Experimental Limitations}
\label{sec:experimental_limitations}

\hs Several limitations should be considered when interpreting the experimental results:

\begin{itemize}
\item \textbf{Repeated holdout evaluation:} The experiments use repeated holdout splits rather than full cross-validation or an external validation set. Although repeated holdout evaluation provides a useful estimate of model behavior across different random splits, it may not be as stable as cross-validation.

\item \textbf{No statistical significance testing:} No statistical significance test is applied in the current evaluation. Therefore, small differences between models should be interpreted descriptively. For example, when two models have very close mean scores and overlapping standard deviations, it is not possible to conclude that one model is statistically better than the other.

\item \textbf{Fixed KFC configurations:} The KFC divergence sets and the number of clusters are fixed in the experiment. They are not selected through a systematic hyperparameter optimization procedure. As a result, the reported KFC results may not represent the best possible performance of the framework.

\item \textbf{Large-dataset sampling strategy:} For large datasets such as Cyber Security Attacks, California Housing, and Wind Turbine, the experiment uses the first 5,000 observations. This sampling strategy makes the experiment computationally manageable, but it may introduce ordering bias if the original dataset is sorted by time, location, or another hidden structure.


\item \textbf{Metric limitation for Wind Turbine:} MAPE is not suitable for the Wind Turbine dataset because the target variable contains values near zero. This can produce undefined, unstable, or inflated percentage errors. Therefore, MAPE is not reported in the final Wind Turbine table, and RMSE and $R^2$ are used as the main metrics for interpretation.
\end{itemize}




\section{Chapter Summary}
\label{sec:chapter_summary_results}

\hs This chapter presented the experimental results of the implemented framework on both classification and regression tasks. The classification experiments evaluated four datasets: Heart Disease, Diabetes, Loan Prediction, and Cyber Security Attacks. The regression experiments evaluated three datasets: Abalone, California Housing Price, and Wind Turbine. All reported results were summarized using the mean and standard deviation across five repeated train--test splits.




The main finding is that the implemented \KFCProcedure{} framework is executable across multiple supervised learning tasks and can be compared consistently with standard baseline models and COBRA-based aggregation models. For classification, \varname{KFCClassifier} obtains the nominally highest reported mean performance on the Heart Disease dataset, but this result should be interpreted cautiously because of the small test set and overlapping standard deviations. On the Diabetes and Loan Prediction datasets, it remains close to strong baselines but does not clearly outperform them. On the Cyber Security Attacks dataset, all models remain near the three-class chance baseline, so the result mainly confirms executability rather than predictive advantage.

For regression, \varname{KFCRegressor} with \varname{GradientCOBRA} completes successfully on all three regression datasets. It obtains a nominal RMSE and $R^2$ advantage on the Abalone dataset, but this advantage is small and within the reported variability. On the California Housing Price and Wind Turbine datasets, the strongest results are obtained by tree-based baseline models or standalone COBRA-based regression models. This suggests that the benefit of \KFCProcedure{} is dataset-dependent and may rely on the presence of meaningful latent cluster structures and suitable local predictive models.

Overall, the experimental evidence supports the usability, modularity, and comparability of the implemented framework. The results also suggest that \KFCProcedure{} may be effective when the data contain useful local structures, but its advantage is not universal across all datasets. Therefore, broader claims of predictive superiority would require additional experiments, including systematic hyperparameter tuning, ablation studies, cross-validation, larger dataset samples, cluster validation analysis, and statistical significance testing.

